\chapter{ConstitutiveState as a Semantic Layer}

\section{Why Another Layer Was Needed}

After the storage cleanup of the previous chapter, the code still had an
architectural ambiguity: \texttt{MaterialState} was doing two jobs at once.

\begin{itemize}
  \item At the low level, it was a generic adapter over a storage policy.
  \item At the constitutive-site level, it was implicitly standing in for the
        local state of a material point or section point.
\end{itemize}

Those roles are related, but they are not identical.  A storage adapter is an
implementation device.  A constitutive state is a modelling concept.

To make that distinction explicit without breaking the existing kernels, this
phase introduces \texttt{ConstitutiveState<StateStoragePolicy, ...>} as a thin,
fully inline semantic wrapper over \texttt{MaterialState}.

\section{Layered View}

The state path is now:

\begin{center}
\texttt{HistoryStorage / CircularHistoryStorage / ElasticState / TrialCommittedState}
$\rightarrow$
\texttt{MaterialState}
$\rightarrow$
\texttt{ConstitutiveState}
$\rightarrow$
\texttt{MaterialInstance}
\end{center}

The meaning of each layer is now cleaner:

\begin{itemize}
  \item \texttt{HistoryStorage} and its variants define how samples are
        physically stored;
  \item \texttt{MaterialState} adapts that storage to a uniform compile-time
        contract;
  \item \texttt{ConstitutiveState} says that the stored quantity is the local
        state of a constitutive site;
  \item \texttt{MaterialInstance} binds that local state to a constitutive law.
\end{itemize}

\section{What ConstitutiveState Adds}

\texttt{ConstitutiveState} does not add runtime overhead or dynamic dispatch.
Its value is semantic and architectural:

\begin{itemize}
  \item it gives \texttt{MaterialInstance} an explicit
        \texttt{ConstitutiveStateT} alias;
  \item it exposes \texttt{constitutive\_state()} as a first-class accessor;
  \item it lets the code and the documentation talk about ``the constitutive
        state of the site'' instead of ``the storage policy wrapper'';
  \item it prepares the next step, where the constitutive law should stop
        owning all irreversible memory internally.
\end{itemize}

In other words, the wrapper is intentionally thin, but its name matters.
Architecturally, names are part of the model.

\section{Effect on MaterialInstance}

\texttt{MaterialInstance<Relation, StatePolicy>} no longer stores
\texttt{MaterialState<StatePolicy, KinematicT>} directly.  It now stores

\begin{center}
\texttt{ConstitutiveState<StatePolicy, KinematicT>}
\end{center}

and exposes it through:

\begin{itemize}
  \item \texttt{current\_state()} as the legacy direct view;
  \item \texttt{constitutive\_state()} as the new semantic accessor.
\end{itemize}

This means user code and future algorithms can progressively move from
``give me the current strain-like value'' toward
``give me the constitutive-state object attached to this site''.

\section{Why This Matters for Future Nonlinear Algorithms}

The next architectural step should separate:

\begin{enumerate}
  \item the constitutive law;
  \item the constitutive state;
  \item the integration algorithm.
\end{enumerate}

That separation is difficult if the local state has no identity of its own.
By introducing \texttt{ConstitutiveState}, the code now has a clear place where
future information can migrate:

\begin{itemize}
  \item committed kinematic state;
  \item trial state;
  \item algorithmic internal variables;
  \item nonlocal history variables;
  \item fracture or damage regularisation data;
  \item time-integration metadata for dynamic constitutive updates.
\end{itemize}

This does not mean all of that should be stored in one monolithic object.  It
means the architecture now has a semantically correct anchor for those future
choices.

\section{Static Polymorphism Is Preserved}

The introduction of \texttt{ConstitutiveState} preserves the compile-time
character of the constitutive hot path:

\begin{itemize}
  \item all storage-policy choices remain template parameters;
  \item no additional virtual calls are introduced;
  \item circular-history windows still use fixed-size arrays and
        compile-time capacity;
  \item the relation type remains statically known inside homogeneous loops.
\end{itemize}

This is important for fiber sections and any future continuum crack/fracture
model where thousands or millions of local constitutive sites may be updated.

\section{What This Refactor Still Does Not Do}

Even with \texttt{ConstitutiveState} in place, the local nonlinear laws in the
current code still keep their irreversible memory internally.  Therefore:

\begin{itemize}
  \item \texttt{ConstitutiveState} is not yet the sole owner of local history;
  \item \texttt{ConstitutiveIntegrator} still does not evolve an explicit state
        object passed by argument;
  \item the design is cleaner, but not yet fully factored.
\end{itemize}

This is acceptable at this stage.  The refactor is intentionally incremental:
first make the ontology correct, then move behaviour.

\section{Immediate Consequence}

From this point onward, the recommended interpretation is:

\begin{itemize}
  \item \texttt{MaterialState}: low-level storage adapter;
  \item \texttt{ConstitutiveState}: local state object of a constitutive site;
  \item \texttt{MaterialInstance}: site-level binding of relation + state;
  \item \texttt{Material<Policy>}: erased owning constitutive handle;
  \item \texttt{MaterialPoint} / \texttt{MaterialSection}: continuum and
        structural constitutive sites bound to geometric integration sites.
\end{itemize}

This vocabulary is much closer to the real mathematical model and therefore a
better base for the next refactor stage.
